This study proposed a dynamic rock--TBM spatial relation graph for component-level geological anomaly interpretation during TBM excavation. The graph represents TSP-derived rock voxels and component-labelled TBM surface nodes as heterogeneous spatial entities, constructs candidate rock--TBM relations under active-zone, distance, surface-normal compatibility and excavation-state constraints, and computes component-level geological anomaly indices from the resulting edge-weighted relations.

The revised evaluation is framed as event-centred representation validation using a confirmed shield-jamming event in the Boshulaling Tunnel. The graph places TSP-derived anomalies, TBM components and monitoring responses in a common excavation coordinate system; reconstructs how low-velocity rock voxels become spatially associated with shield components during event development; compares component-level indices with alternative summaries and ablated graph settings; and traces the interpretation back to specific rock voxels, TBM surface nodes and candidate edges.

The primary value of the proposed graph is to make spatial relations between anomalous geological units and moving TBM components explicit, component-level and traceable. It should be understood as a queryable spatial relation representation for event interpretation and future model integration, rather than as a standalone calibrated failure-prediction or contact-force estimation model.
